\begin{proofbox}
	\textbf{\textcolor{CProof}{思路提示}}

	先验证水平变换如何作用于同一个点，再将纵坐标乘以带符号的 $A$。设 $A\neq0$，$\omega>0$，$\varphi\in\mathbb R$。

	\medskip
	\textbf{\textcolor{CProof}{两条函数变换链}}
	\begin{description}[style=nextline, leftmargin=2.8em, labelsep=0.5em, itemsep=0.55em]
		\item[\textbf{先移后缩：}] 向左平移有符号距离 $\varphi$，再将横坐标乘以 $1/\omega$，最后将纵坐标乘以 $A$：
			\[
				\sin x\ \longrightarrow\ \sin(x+\varphi)
				\ \longrightarrow\ \sin(\omega x+\varphi)
				\ \longrightarrow\ A\sin(\omega x+\varphi).
			\]
		\item[\textbf{先缩后移：}] 先将横坐标乘以 $1/\omega$，再向左平移有符号距离 $\varphi/\omega$，最后将纵坐标乘以 $A$：
			\[
				\sin x\ \longrightarrow\ \sin(\omega x)
				\ \longrightarrow\ \sin\!\Big[\omega\Big(x+\frac{\varphi}{\omega}\Big)\Big]
				\ \longrightarrow\ A\sin(\omega x+\varphi).
			\]
	\end{description}

	\medskip
	\textbf{\textcolor{CProof}{点坐标验证}}

	取基准曲线上的任意点 $P=(t,\sin t)$。第一条路径把它送到
	\[
		(t,\sin t)\longrightarrow(t-\varphi,\sin t)
		\longrightarrow\left(\frac{t-\varphi}{\omega},\sin t\right)
		\longrightarrow\left(\frac{t-\varphi}{\omega},A\sin t\right).
	\]
	第二条路径把同一点送到
	\[
		(t,\sin t)\longrightarrow\left(\frac{t}{\omega},\sin t\right)
		\longrightarrow\left(\frac{t}{\omega}-\frac{\varphi}{\omega},\sin t\right)
		\longrightarrow\left(\frac{t-\varphi}{\omega},A\sin t\right).
	\]
	两条路径的最终点相同。令其横坐标为 $x=(t-\varphi)/\omega$，则 $t=\omega x+\varphi$，纵坐标恰为 $A\sin(\omega x+\varphi)$。由于 $t$ 可取任意实数、$\omega>0$，两条路径都得到整幅目标图像。

	\medskip
	\textbf{\textcolor{CProof}{符号与顺序}}

	若 $A<0$，最后一步可分解为纵坐标乘以 $|A|$，再关于 $x$ 轴对称；振幅仍为 $|A|$。横向点变换满足
	\[
		\frac{t-\varphi}{\omega}=
		\frac{t}{\omega}-\frac{\varphi}{\omega},
	\]
	因此调整平移量后两条路径相同。若先伸缩后仍平移 $\varphi$，一般不能得到同一图像；这不是平移与伸缩的普通交换律。
\end{proofbox}
